\documentclass[11pt]{article}

\usepackage[margin=1in]{geometry}
\usepackage{amsmath,amssymb}
\usepackage{booktabs}
\usepackage{hyperref}
\usepackage{microtype}

\newcommand{\R}{\mathbb{R}}
\newcommand{\Xs}{\mathcal{X}}
\newcommand{\GP}{\mathcal{GP}}
\newcommand{\Norm}{\mathcal{N}}
\newcommand{\Data}{\mathcal{D}}

\title{Separating Two Things Composite Modeling Provides:\\
Knowing the Outer Function from Observing the Intermediates}
\author{}
\date{}

\begin{document}
\maketitle

\section{Where this sits}

Our benchmark study compared two surrogates on problems of the composite form
$f(x)=g(h(x))$, where $h$ is an expensive black box and $g$ is a known, cheap
function. \emph{Direct} fits a Gaussian process to $f$; \emph{Composite} fits
Gaussian processes to $h$ and pushes their posterior samples through $g$.
Composite improved most comparisons, but by amounts ranging from $121\%$ to
nothing, and the follow-up experiments are designed to find out what separates
those cases.

This note addresses a confound that runs through all of them. Composite gives
the surrogate \textbf{two different things at the same time}:

\begin{enumerate}
\item the \textbf{knowledge of $g$} --- the surrogate is told the algebraic form
of the final objective, so it never has to learn that part from data;
\item the \textbf{observations of $h$} --- each expensive evaluation returns
$k$ numbers instead of one.
\end{enumerate}

Every Direct--Composite comparison we have run, and every experiment currently
planned, changes both at once. So when Composite wins by $121\%$ on DTLZ2 we
cannot say which of the two is responsible, and when it wins by nothing on
RGB we cannot say which one failed to pay off.

The fix is to add a third surrogate that receives exactly one of the two.

\section{Three arms}

Fix a single objective $f(x)=g(h(x))$ with $h:\Xs\to\R^k$ unknown and
$g:\R^k\to\R$ known. Write $\Data_n=\{(x_i,\cdot)\}_{i\le n}$ for the data after
$n$ evaluations.

\begin{center}
\begin{tabular}{llll}
\toprule
Arm & Observes & Knows $g$ & Surrogate \\
\midrule
1. Direct    & $y_i=f(x_i)$ & no  & $f\sim\GP$ fit to $y$ \\
2. Composite & $h(x_i)$     & yes & $h_j\sim\GP$ fit to $h_j$, then apply $g$ \\
3. Latent    & $y_i=f(x_i)$ & yes & $h_j\sim\GP$ \emph{a priori}, $h$ never observed \\
\bottomrule
\end{tabular}
\end{center}

Arm 3 is the new one. It has the same prior belief about $h$ that Composite
has, and it knows $g$ exactly, but it only ever sees the scalar $y$. Its
posterior over the latent values $\mathbf{h}(X)=(h(x_1),\dots,h(x_n))$ is
just Bayes' rule,
\begin{equation}
p\bigl(\mathbf{h}(X)\mid \mathbf{y}\bigr)\;\propto\;
\underbrace{\prod_{i=1}^{n}\Norm\!\bigl(y_i;\,g(h(x_i)),\,\sigma^2\bigr)}_{\text{likelihood}}
\;\cdot\;
\underbrace{\prod_{j=1}^{k}\Norm\!\bigl(h_j(X);\,m_j,\,v_jK\bigr)}_{\text{GP prior}} ,
\label{eq:posterior}
\end{equation}
and prediction at a new point is exact given the latent values, because
$h(x')\mid \mathbf{h}(X)$ is an ordinary GP conditional. The arm-3 predictive is
therefore a \emph{mixture of Gaussian processes}, one component per posterior
sample of $\mathbf{h}(X)$.

Because arms 2 and 3 use the \textbf{identical} prior on $h$ --- the true kernel
of the generator that produced the test functions --- the only difference
between them is what they observe. The two comparisons then read cleanly:
\begin{align}
\text{Arm 3} - \text{Arm 1} &\;=\; \text{the value of \emph{knowing} } g, \\
\text{Arm 2} - \text{Arm 3} &\;=\; \text{the value of \emph{observing} } h .
\end{align}

\section{Why the decomposition should be informative}

The second quantity has a clean characterization: observing $h$ can only help
to the extent that $y=g(h)$ fails to determine $h$. That makes the
\emph{invertibility of $g$} the governing property, and it gives three regimes
we can reason about in advance.

\subsection{Invertible $g$: observing $h$ is worth nothing}

If $g$ is injective and the observations are noise-free, then
$h(x_i)=g^{-1}(y_i)$ exactly. Arm 3 recovers the intermediates from the
objective alone, so arms 2 and 3 condition on identical information and are the
\emph{same model}. Any advantage Composite shows over Direct on such a problem
is attributable entirely to knowing $g$.

The doc's exponential warping family, $g(z)=\exp(\beta z)$, is of this type.
Knowing $g$ is still worth something --- it tells the surrogate that $y>0$ and
how the scale is warped --- but the extra numbers are redundant.

\subsection{Linear $g$: knowing $g$ is worth nothing}

If $g(\mathbf{y})=\sum_j c_j y_j$ and the $h_j$ are independent GPs, then
$g(h(\cdot))$ is \emph{itself} a GP, with mean $\sum_j c_jm_j$ and kernel
$\bigl(\sum_j c_j^2 v_j\bigr)k(\cdot,\cdot)$. Arm 3 therefore collapses to an
ordinary GP fitted to $y$ --- which is arm 1 with a particular kernel. Up to
hyperparameter estimation, knowing $g$ buys nothing at all.

This is the sharpest prediction in the design, and it matters for the paper:
Truss and RCM40 both have linear outer maps, and both produced our weakest
results. If the prediction holds, the explanation is not merely that ``linear
maps are easy'' but something more specific --- on a linear map there is no
knowledge to exploit, so whatever gain Composite shows must come from the extra
observations alone.

\subsection{Non-invertible $g$: both channels are live}

Take the smallest interesting case, $g(z)=z^2$, with a noise-free observation
$y=g(h(x))$. Then $h(x)$ is known only up to sign,
\begin{equation}
p\bigl(h(x')\mid g(h(x))=y\bigr)
\;=\;p\bigl(h(x')\;\bigm|\;h(x)\in\{+\sqrt{y},\,-\sqrt{y}\}\bigr),
\label{eq:square}
\end{equation}
so the posterior is a two-component mixture for one observation, and a mixture
over $2^n$ sign patterns for $n$ of them, weighted by how plausible each pattern
is under the GP prior. Observing $h$ removes that ambiguity outright, which is
what arm 2 measures.

But knowing $g$ is worth something even without any observation of $h$. The
range of $g$ is $[0,\infty)$, so arm 3 assigns
\begin{equation}
P\bigl(y(x')=-1\mid\Data\bigr)=P\bigl(h(x')^2=-1\mid\Data\bigr)=0 ,
\end{equation}
whereas a Direct GP on $y$ --- a Gaussian --- puts positive probability on
negative values it can never observe. More generally, knowing $g$ imposes the
support, the shape, and any kinks or saturation of the outer map, none of which
a stationary GP on $y$ can represent.

In general the information destroyed by $g$ is the \emph{fiber} $g^{-1}(y)$: a
finite set for $z^2$ or $\cos$, a half-line for a clipped or rectified map, a
circle for a Euclidean norm, a hyperbola for a product. This is the structure
the synthetic maps in Section~\ref{sec:functions} are chosen to vary.

\subsection{Summary of predictions}

\begin{center}
\begin{tabular}{llcc}
\toprule
$g$ & Type & Arm 3 $-$ Arm 1 & Arm 2 $-$ Arm 3 \\
& & (knowing $g$) & (observing $h$) \\
\midrule
$z$                    & identity    & $0$      & $0$ \\
$\exp(\beta z)$, $\sinh(\beta z)$ & invertible & $+$ & $0$ by construction \\
$(z-c)^2$              & lossy, tunable & $+$   & $+$ at $c=0$, $\to 0$ as $c$ grows \\
$|z|$                  & lossy       & $+$      & $+$ \\
$y_1+y_2$              & linear      & $\approx 0$ & $+$ \\
norm, max, product     & lossy       & $+$      & $+$ \\
\bottomrule
\end{tabular}
\end{center}

\section{Test functions}
\label{sec:functions}

\subsection{Design principle}

The intermediates must not become a confound. Every primitive $h_j$ is
therefore a fixed draw from the same RBF-GP family used by the surrogate,
generated once with random Fourier features,
\begin{equation}
h(x)=\sqrt{\tfrac{2}{M}}\sum_{m=1}^{M}a_m\cos(\omega_m^\top x+b_m),
\qquad a_m\sim\Norm(0,1),\;\;b_m\sim U(0,2\pi),\;\;\omega_m\sim\Norm(0,\ell^{-2}I),
\end{equation}
with $M=1024$ and $\ell=0.5$, which has prior $\GP(0,k_{\mathrm{RBF}})$. The draw
is standardized on a fixed Sobol reference set, so the implied prior on the
standardized function is $\GP(-\mu/s,\,k_{\mathrm{RBF}}/s^2)$; those constants are
what arms 2 and 3 use as their known prior. The same draw seeds are reused for
every condition, so switching $g$ never changes the underlying functions.

Standardization matters for a reason beyond scale: it puts the draw's mean at
zero, so $h$ takes \emph{both signs}. Without that, $z^2$ would be injective over
the region actually visited and would lose nothing.

Each final objective is normalized to roughly $[0,1]$ using the $1$st and $99$th
percentiles of $g(h(\cdot))$ on a fixed reference set, so that different outer
maps do not differ merely in scale. The normalization is affine and known, so it
is part of $g$ for all three arms.

\subsection{The maps}

\paragraph{The sweep.} $g_c(z)=(z-c)^2$ for $c\in\{0,0.5,1,2,3\}$ holds the
functional form fixed and varies only how much ambiguity there is. At $c=0$ the
sign of $h$ is lost wherever the draw crosses zero; at $c=3$ the shifted argument
almost never changes sign over the domain, so $g$ is effectively injective. This
is the cleanest available knob on ``how much $g$ discards.''

\paragraph{Matched pairs.} Each pair changes one property and holds the rest:
\begin{center}
\begin{tabular}{lll}
\toprule
Pair & Shared & Isolates \\
\midrule
$z^2$ vs.\ $|z|$ & both two-to-one & smoothness at the optimum \\
$\exp(\beta z)$ vs.\ $\sinh(\beta z)$ & both injective, both smooth & whether a range constraint helps \\
$\cos(\omega z)$, $\omega\in\{1,2,4\}$ & smooth & \emph{how many} preimages \\
$\max(z,0)$, $\mathrm{clip}(cz,-1,1)$ & --- & censoring and saturation \\
\bottomrule
\end{tabular}
\end{center}

\paragraph{Benchmark analogs.} Two intermediates into one objective, each
mirroring a real benchmark: $y_1+y_2$ (Truss, RCM40), $y_1y_2$ (DTLZ2),
$\sqrt{y_1^2+y_2^2}$ (Color), $\max(y_1,y_2)$ (Lake), and
$u_1/(0.05+u_1+\lambda u_2)$ with $u=\mathrm{sigmoid}(y)$ (RGB). These let the
synthetic conclusions be carried back to the benchmark results.

\section{Inference for the latent arm}

Arm 3 is the only one that needs more than a closed form. Three cases:

\paragraph{Closed form where possible.} For the identity and injective maps,
$h=g^{-1}(y)$ exactly, so arm 3 is implemented as arm 2. For a linear map the
posterior \eqref{eq:posterior} is Gaussian and is available analytically. Both
shortcuts are validated against the general sampler in Section~\ref{sec:validation}.

\paragraph{General case.} Otherwise we sample \eqref{eq:posterior} over the
$n\times k$ latent values, which is at most $90$ numbers at our budget. Two move
types alternate.

\emph{Elliptical slice sampling} \cite{murray2010} is the standard sampler for a
GP prior with an arbitrary likelihood: it proposes
$\mathbf{h}\cos\theta+\boldsymbol{\nu}\sin\theta$ with
$\boldsymbol{\nu}$ a prior draw, and shrinks $\theta$ until the likelihood
threshold is met. It needs no gradients, which matters because $|z|$, $\max$ and
the clipped map have kinks or flat regions, and no tuning.

\emph{Level-set moves} do the work that elliptical slice sampling cannot. When
$\sigma$ is small the likelihood is sharply peaked on the set
$\{g(h(x_i))=y_i\}$, and slice proposals must shrink to roughly
$\sigma/\mathrm{sd}(h)$ to stay on it, so movement \emph{along} that set is a slow
random walk --- precisely along the directions the experiment is about. But
those directions are known in closed form for each $g$. We therefore propose
$h(x_i)\mapsto h'(x_i)$ with
\begin{equation}
g\bigl(h'(x_i)\bigr)=g\bigl(h(x_i)\bigr) \text{ exactly},
\end{equation}
from a symmetric proposal, so the likelihood ratio is identically one and the
Metropolis--Hastings acceptance probability reduces to the GP prior ratio alone.
With $Q=K^{-1}$ and $r_j=h_j(X)-m_j$, changing coordinate $i$ of latent $j$ from
$\alpha$ to $\beta$ gives
\begin{equation}
\Delta\log p=-\frac{1}{2v_j}\Bigl[Q_{ii}\bigl(\beta^2-\alpha^2\bigr)
+2(\beta-\alpha)\bigl((Qr_j)_i-Q_{ii}\alpha\bigr)\Bigr].
\end{equation}
The moves are the symmetries of each map: reflection $h\mapsto 2c-h$ for
$(z-c)^2$ and $|z|$; reflection through the nearest multiple of $\pi/\omega$ plus
translation by the period for $\cos(\omega z)$; a random walk confined to the
flat region for $\max(z,0)$ and the clipped map; rotation of $(h_1,h_2)$ for the
norm; negation and reciprocal rescaling for the product; swapping the two
coordinates and moving the smaller one for the max; and the shift
$(h_1+\delta,h_2-\delta)$ for the sum. Only the ratio map has no convenient
exact symmetry, and is sampled by elliptical slice sampling alone.

\paragraph{Practical details.} Burn-in anneals $\sigma$ geometrically down to its
target so that chains started from the prior can find the data; later fits
warm-start from the previous chain, with the new point initialized from its GP
conditional. Several chains are run and diagnostics are computed on the
predictive of $f$, which is invariant to the symmetries of $h$ that leave $f$
unchanged --- so a chain that has globally flipped sign under an even $g$ does
not register as non-convergence.

\section{What is measured}

\paragraph{Information destroyed by $g$.} The average posterior variance of the
intermediates \emph{at the points already evaluated},
\begin{equation}
U_{\text{lost}}=\frac{1}{nk}\sum_{i=1}^{n}\sum_{j=1}^{k}
\operatorname{Var}\bigl(h_j(x_i)\mid \mathbf{y}\bigr),
\end{equation}
is zero when $g$ is injective and grows with the size of the fiber. It is a
measured version of the quantity the whole design is organized around, and the
central plot is the ``observing $h$'' gain against it.

\paragraph{Optimization performance.} Each arm reports simple regret against a
reference minimum; comparisons use the final $\log_{10}$ regret, paired over
seeds, since all three arms share an initial design and an acquisition
procedure.

\paragraph{A confound to control.} Arms 2 and 3 are given the generator's true
kernel, whereas arm 1 must estimate its hyperparameters by maximum likelihood,
because for a general $g$ there is no true kernel for $f$. The ``knowing $g$''
comparison therefore mixes the value of $g$ with the value of a correct prior.
The identity map measures the second effect alone --- knowing $g$ is worth
exactly nothing there --- and so serves as the baseline to subtract. A
robustness run with maximum-likelihood hyperparameters in all three arms
resolves it directly.

\section{Protocol}

Single objective, $d=6$, $\ell=0.5$; $n_0=5$ Sobol points and $40$ adaptive
evaluations, matching Experiment~1; $10$ paired seeds. All three arms share the
initial design, the acquisition function (Monte Carlo LogEI), the candidate pool
and the evaluation budget, so the surrogate is the only thing that differs.

\section{Validation}
\label{sec:validation}

The sampler is checked against the two cases whose answers are known
independently before being trusted anywhere else.

\paragraph{Linear map, against the exact posterior.} On $y_1+y_2$ with $n=20$,
the sampled posterior matches the analytic Gaussian: the error in the posterior
mean is a median of $0.065$ posterior standard deviations (max $0.249$), and the
ratio of sampled to exact standard deviation has median $0.973$ and range
$[0.85,1.10]$.

\paragraph{Injective map, for calibration.} On $\exp(z)$ the true $h$ falls
within two posterior standard deviations at $95\%$ of points. The largest raw
error, $0.62$, occurs at $h=-1.96$, where $\exp$ is flat enough that its slope
falls below the likelihood noise. This is correct behavior and a substantive
caveat: \textbf{with observation noise, an injective $g$ still destroys
information wherever it is flat}, so ``arm 2 equals arm 3 for injective $g$''
holds exactly only in the noise-free limit.

\paragraph{The sweep behaves as designed.}

\begin{center}
\begin{tabular}{lrrr}
\toprule
Map & $U_{\text{lost}}$ & Chain disagreement & Move acceptance \\
\midrule
$(z-0)^2$   & $0.877$ & $0.029$ & $0.67$ \\
$(z-0.5)^2$ & $1.072$ & $0.052$ & $0.52$ \\
$(z-1)^2$   & $0.443$ & $0.047$ & $0.32$ \\
$(z-2)^2$   & $0.139$ & $0.036$ & $0.07$ \\
$(z-3)^2$   & $0.003$ & $0.018$ & $0.00$ \\
$|z|$       & $0.882$ & $0.027$ & $0.66$ \\
\bottomrule
\end{tabular}
\end{center}

The information destroyed falls by more than two orders of magnitude across the
sweep, and the level-set acceptance rate falls with it --- at $c=3$ the
reflected proposal is so implausible under the prior that it is never accepted,
which is the correct behavior when there is no ambiguity left to explore. Chain
disagreement is in normalized objective units, where the objective spans about
one.

\section{Limitations}

\begin{itemize}
\item \textbf{Noise makes injectivity a matter of degree.} As measured above,
an injective $g$ still loses information where its slope is small relative to
$\sigma$. The clean statement holds only in the noise-free limit.
\item \textbf{Continuous fibers mix more slowly than discrete ones.} Reflections
for $z^2$ are single moves that jump between modes; following the circle fiber of
a norm is a single-site random walk over a smooth field, which mixes more slowly
when observed points are strongly coupled. The ratio map has no exact moves at
all and relies on elliptical slice sampling.
\item \textbf{The reference minimum is estimated}, from a large Sobol set
refined by gradient ascent, so regret is measured against a very good estimate
rather than a proven optimum.
\item \textbf{Single objective.} The decomposition is cleanest with one
objective; whether the same accounting holds once a hypervolume acquisition is
involved is a separate question, and the natural next step once these results
are in.
\end{itemize}

\begin{thebibliography}{9}
\bibitem{murray2010}
I.~Murray, R.~P. Adams, and D.~J.~C. MacKay.
Elliptical slice sampling.
In \emph{Proceedings of the 13th International Conference on Artificial
Intelligence and Statistics}, PMLR 9:541--548, 2010.
\end{thebibliography}

\end{document}
